\documentclass[11pt]{article}
\usepackage[margin=0.8in,landscape]{geometry}
\usepackage{booktabs}
\usepackage{amsmath,amssymb}
\usepackage{caption}
\usepackage{threeparttable}
\usepackage[T1]{fontenc}
\usepackage{lmodern}
\captionsetup[table]{labelfont=bf, font=small, skip=6pt}
\begin{document}
\begin{table}[htbp]
\centering
\footnotesize
\begin{threeparttable}
\caption{Leverage dynamics around the GFC and pre-crisis concentration, by region and country}
\label{tab:region_gfc}
\setlength{\tabcolsep}{4pt}
\begin{tabular}{lrrrrrrrrrrr}
\toprule
 & & \multicolumn{2}{c}{Top quantile} & \multicolumn{2}{c}{Median} & \multicolumn{6}{c}{Structure in 2006} \\
\cmidrule(lr){3-4}\cmidrule(lr){5-6}\cmidrule(lr){7-12}
 & Banks & Build-up & Delever. & Build-up & Delever. & CR5 $A$ & CR5 $E$ & Top10\%-lev $A$ sh. & $\rho$(size, lev) & Lev. D10/D1 & Median lev. \\
\midrule
\multicolumn{12}{l}{\textit{Regions (top = p95)}} \\
N. America & 203 & 3.4 & -19.5 & -0.8 & -14.4 & 51.3 & 45.1 & 29.8 & 0.15 & 1.65 & 10.8 \\
Europe & 289 & 21.4 & -25.5 & 5.1 & -15.8 & 27.2 & 27.5 & 19.8 & 0.36 & 2.33 & 17.7 \\
Japan \& Dev. AP & 118 & 36.2 & -43.4 & 13.5 & -14.8 & 48.7 & 43.9 & 5.5 & 0.04 & 1.78 & 16.5 \\
China & 20 & 16.3 & -35.7 & -2.3 & -14.1 & 89.9 & 93.2 & 4.6 & -0.17 & 0.75 & 18.8 \\
Emerging Asia & 67 & -11.7 & -10.4 & -2.1 & -9.6 & 39.6 & 34.7 & 7.9 & 0.15 & 1.28 & 12.8 \\
Latin America & 55 & 4.9 & -8.2 & 2.7 & 6.0 & 53.4 & 55.4 & 10.9 & 0.38 & 1.62 & 10.5 \\
M. East \& Africa & 68 & -15.2 & -12.0 & 2.8 & -18.4 & 34.3 & 28.1 & 11.0 & 0.37 & 2.24 & 10.1 \\
Emerging Eur. \& CIS & 24 & -10.7 & 13.1 & 7.0 & -5.6 & 75.7 & 72.9 & 10.3 & 0.05 & 1.31 & 8.5 \\
\multicolumn{12}{l}{\textit{Europe split (top = p90)}} \\
Europe core & 205 & 6.8 & -20.2 & 7.2 & -19.0 & 40.0 & 35.3 & 22.1 & 0.44 & 2.48 & 18.2 \\
Europe periphery & 49 & 12.0 & 22.9 & 0.9 & 0.9 & 58.2 & 57.0 & 8.5 & 0.20 & 1.90 & 15.5 \\
UK & 21 & 34.3 & -45.7 & 20.8 & -34.5 & 86.5 & 89.4 & 31.9 & 0.37 & 1.92 & 20.9 \\
\multicolumn{12}{l}{\textit{Countries (top = p90)}} \\
United States & 184 & 2.5 & -20.1 & -1.7 & -14.6 & 61.1 & 50.1 & 41.8 & 0.06 & 1.31 & 10.5 \\
Germany & 85 & 23.8 & -44.7 & 0.2 & -22.3 & 55.7 & 46.1 & 33.8 & 0.57 & 2.43 & 21.3 \\
Japan & 64 & 35.4 & -33.7 & 26.8 & -24.3 & 69.0 & 67.2 & 5.3 & -0.16 & 1.06 & 17.5 \\
Switzerland & 30 & 4.3 & 6.4 & 3.1 & 5.2 & 81.8 & 60.7 & 72.8 & 0.30 & 3.02 & 12.3 \\
Italy & 24 & 59.4 & 32.0 & -7.2 & 7.5 & 80.9 & 79.1 & 4.8 & 0.16 & 1.41 & 14.6 \\
United Kingdom & 21 & 34.3 & -45.7 & 20.8 & -34.5 & 86.5 & 89.4 & 31.9 & 0.37 & 1.92 & 20.9 \\
Canada & 19 & 1.1 & -19.0 & -1.3 & -16.7 & 82.8 & 80.8 & 1.3 & 0.11 & 0.94 & 21.2 \\
France & 19 & 15.9 & -6.6 & -12.7 & 9.3 & 95.8 & 91.8 & 1.1 & 0.66 & 12.37 & 14.2 \\
Taiwan & 19 & 0.5 & -11.4 & 3.2 & -9.9 & 53.0 & 48.8 & 20.4 & -0.18 & 1.43 & 16.5 \\
China & 18 & 12.0 & -26.2 & -4.4 & -16.5 & 90.1 & 93.5 & 4.6 & -0.30 & 0.75 & 19.8 \\
Austria & 18 & 2.9 & -1.2 & 15.9 & -9.4 & 75.7 & 82.3 & 4.1 & -0.01 & 1.02 & 16.6 \\
\bottomrule
\end{tabular}
\begin{tablenotes}[flushleft]
\footnotesize
\item \textit{Notes:} BankFocus global panel (banks with total assets $\geq$ USD 10bn in at least one year), constant sample of banks observed in every year 2006--2012 within each group; French mutualist network members excluded; leverage = total assets / total equity, winsorized at 1/99\% by year. Unweighted leverage quantiles indexed to 2006 = 100. Build-up = $\max(I_{2007}, I_{2008}) - 100$; Delever. = $I_{2012} - \max(I_{2007}, I_{2008})$ (index points). Top quantile is p95 for regions, p90 for the Europe split and countries (fewer banks). CR5 $A$ ($E$): share of the five largest banks in total assets (equity), \%. Top10\%-lev $A$ sh.: asset share of the 10\% most levered banks, \%. $\rho$(size, lev): Spearman correlation of assets and leverage. Lev. D10/D1: mean leverage in the top over the bottom asset decile. Median lev.: median leverage. All structure measures in 2006. Europe core = DE, FR, NL, AT, BE, LU, CH and Nordics; periphery = IT, ES, PT, GR, IE.
\end{tablenotes}
\end{threeparttable}
\end{table}

\clearpage

\begin{table}[htbp]
\centering
\footnotesize
\begin{threeparttable}
\caption{Robustness and decomposition: build-up and deleveraging by region and country}
\label{tab:region_gfc_decomp}
\setlength{\tabcolsep}{4pt}
\begin{tabular}{lrrrrrrrrrr}
\toprule
 & \multicolumn{4}{c}{Split windows (index pts)} & \multicolumn{4}{c}{Top-decile banks (log pts)} & \multicolumn{2}{c}{Other banks (log pts)} \\
\cmidrule(lr){2-5}\cmidrule(lr){6-9}\cmidrule(lr){10-11}
 & Banks & Banks & Build-up & Delever. & $\Delta A$ & $\Delta E$ & $\Delta A$ & $\Delta E$ & $\Delta A$ & $\Delta E$ \\
 & 06--08 & 08--12 & top & top & 06--08 & 06--08 & 08--12 & 08--12 & 06--08 & 06--08 \\
\midrule
\multicolumn{11}{l}{\textit{Region}} \\
N. America & 233 & 218 & 16.7 & -18.4 & -2.7 & 31.4 & 23.5 & 40.1 & 27.6 & 25.3 \\
Europe & 317 & 307 & 25.6 & -20.8 & 17.9 & 9.9 & -20.9 & 31.0 & 30.4 & 15.2 \\
Japan \& Dev. AP & 124 & 133 & 39.4 & -26.4 & 19.9 & -29.6 & 22.9 & 85.6 & 19.7 & -7.7 \\
China & 24 & 36 & 10.6 & -27.9 & 69.0 & 62.5 & 92.7 & 146.1 & 45.7 & 46.4 \\
Emerging Asia & 71 & 73 & -8.5 & -9.0 & 22.8 & 55.9 & 60.1 & 84.1 & 21.9 & 20.2 \\
Latin America & 59 & 59 & 3.2 & -6.4 & 8.3 & 48.2 & 74.2 & 79.2 & 50.4 & 27.3 \\
M. East \& Africa & 72 & 75 & -15.2 & -17.2 & 23.7 & 34.7 & 22.8 & 49.3 & 34.4 & 34.1 \\
Emerging Eur. \& CIS & 24 & 30 & -10.7 & 0.6 & 31.9 & 58.8 & 41.4 & 61.7 & 54.3 & 52.3 \\
\multicolumn{11}{l}{\textit{Europe split}} \\
Europe core & 221 & 218 & 10.3 & -21.2 & 23.1 & 11.8 & -19.5 & 28.7 & 23.4 & 7.8 \\
Europe periphery & 58 & 52 & 12.4 & 26.2 & 14.4 & -50.2 & -78.9 & 47.7 & 31.0 & 36.4 \\
UK & 22 & 21 & 34.6 & -34.1 & 33.6 & 26.5 & -28.0 & 38.2 & 40.9 & 2.0 \\
\multicolumn{11}{l}{\textit{Country}} \\
United States & 213 & 195 & 5.8 & -18.9 & -10.9 & 27.0 & 1.7 & 24.7 & 36.6 & 26.8 \\
Germany & 85 & 86 & 23.8 & -39.0 & 32.1 & -2.9 & -16.8 & 46.5 & 13.2 & 0.9 \\
Japan & 68 & 72 & 40.1 & -22.4 & 22.9 & -25.8 & 16.9 & 79.2 & 20.3 & -16.3 \\
Switzerland & 32 & 31 & -5.7 & 7.4 & 9.6 & 69.9 & 0.2 & 8.5 & 26.8 & 23.9 \\
Italy & 27 & 25 & 38.6 & 28.7 & 65.9 & 37.2 & 24.8 & 36.8 & 41.9 & 52.2 \\
United Kingdom & 22 & 21 & 34.6 & -34.1 & 33.6 & 26.5 & -28.0 & 38.2 & 40.9 & 2.0 \\
Canada & 20 & 23 & -15.8 & -17.3 & 94.9 & 56.4 & -80.4 & 43.7 & 19.0 & 19.7 \\
France & 19 & 20 & 15.9 & -9.0 & 32.5 & 16.0 & 37.1 & 25.9 & 32.2 & 15.3 \\
Taiwan & 20 & 23 & -5.3 & -8.3 & 5.7 & -104.8 & 30.9 & 134.7 & 3.1 & -5.1 \\
China & 22 & 33 & 11.8 & -23.8 & 69.0 & 62.5 & 92.7 & 146.1 & 45.8 & 46.5 \\
Austria & 23 & 19 & 13.8 & -6.3 & 25.5 & 108.7 & -47.8 & -54.0 & 28.2 & 19.3 \\
\bottomrule
\end{tabular}
\begin{tablenotes}[flushleft]
\footnotesize
\item \textit{Notes:} Split windows: build-up uses banks observed in every year 2006--2008 (so banks that failed or were absorbed after 2008 are kept), $\max(I_{2007}, I_{2008}) - 100$ with 2006 = 100; deleveraging uses banks observed in every year 2008--2012, $I_{2012} - 100$ with 2008 = 100. Top quantile = p95 for regions, p90 for the Europe split and countries; groups with fewer than 15 banks are not reported. Decomposition: within the 2006--2012 constant sample, banks in the top decile of 2006 leverage vs all others; $\Delta A$, $\Delta E$ are $100\times$ log changes of the group's summed total assets and equity (USD at closing-date exchange rates). Leverage rises when $\Delta A > \Delta E$.
\end{tablenotes}
\end{threeparttable}
\end{table}
\end{document}
